\chapter{Justificativa e Histórico}
\label{chp:justificativa}

O \ppccurso da Universidade Federal de Uberlândia (UFU) é uma proposta de curso de graduação tecnológica vinculada à Faculdade de Engenharia Elétrica (FEELT), estruturada em conformidade com o Catálogo Nacional de Cursos Superiores de Tecnologia (CNCST, 4ª edição)~\cite{mec_cncst_2024} e com as Diretrizes Curriculares Nacionais Gerais para a Educação Profissional e Tecnológica~\cite{cne_cp_1_2021}. A proposta parte da experiência consolidada da FEELT na formação em eletricidade, eletrônica, controle e automação — áreas em que a unidade mantém corpo docente qualificado e infraestrutura laboratorial estabelecida — para oferecer uma formação tecnológica de ciclo curto, aplicada e diretamente orientada ao exercício profissional em automação industrial. A criação deste curso foi viabilizada pela contemplação da FEELT pelo Programa Universidades Inovadoras e Sustentáveis~\cite{programa_universidades_inovadoras_sustentaveis} (\autoref{sec:historico}), no âmbito do qual a unidade recebeu novas vagas docentes e de técnicos de laboratório — o mesmo programa que sustentou a reformulação e renomeação do curso de Engenharia de Controle e Automação para Engenharia de Automação, Robótica e Inteligência Artificial, com o qual este curso passa a integrar a Área Básica de Ingresso (ABI, \autoref{chp:abi}; \autoref{sec:distincao_perfis}).

\section{Justificativa}

A expansão acelerada da automação industrial, impulsionada pela convergência entre eletrônica, controle, redes industriais, sistemas ciberfísicos, Internet Industrial das Coisas (IIoT) e inteligência artificial aplicada — o conjunto de transformações associado à Indústria 4.0 —, impôs uma demanda crescente por profissionais de formação tecnológica de ciclo curto, capazes de atuar diretamente no projeto, na implementação, na operação e na manutenção de sistemas automatizados industriais. Ao mesmo tempo, o setor produtivo exige desses profissionais um perfil ético, comunicativo, colaborativo e sensível às realidades socioculturais em que estão inseridos.

O \ppccurso responde também ao chamado das políticas públicas e institucionais que estabelecem novos compromissos para os cursos de graduação, com destaque para a curricularização da extensão — historicamente associada à meta 12.7 do Plano Nacional de Educação 2014--2024~\cite{mec_pne_2014}, hoje regida pelo Plano Nacional de Educação 2026--2036~\cite{mec_pne_2026} — e, no âmbito da UFU, pela Resolução CONGRAD nº 177/2026~\cite{ufu_congrad_177_2026} (organização modular das Atividades Curriculares de Extensão detalhada em \autoref{sec:ace}). A Universidade Federal de Uberlândia, por meio do seu Plano Institucional de Desenvolvimento e Expansão (PIDE)~\cite{ufu_consun_31_2022}, orienta-se por essas diretrizes ao estimular a inovação pedagógica, a interdisciplinaridade, a equidade e a responsabilidade social, e a formação orientada por competências prevista pela própria Resolução CNE/CP nº 1/2021~\cite{cne_cp_1_2021} para a Educação Profissional e Tecnológica.

A proposta também parte de um olhar atento para o contexto regional do Triângulo Mineiro e do interior do Brasil. A vocação tecnológica e industrial da região, os arranjos produtivos locais e a presença de indústrias com processos automatizados demandam uma formação versátil, aplicada e conectada às necessidades concretas do setor produtivo local, nacional e global — características centrais de um curso superior de tecnologia.

\subsection{Justificativa da Denominação Regulatória e da Ênfase Curricular}

A denominação regulatória do curso — \ppccurso — adota literalmente a nomenclatura do Catálogo Nacional dos Cursos Superiores de Tecnologia (CNCST, 4ª edição, MEC/SETEC, 2024)~\cite{mec_cncst_2024}, eixo tecnológico Controle e Processos Industriais, em conformidade com a Resolução CNE/CP nº 1/2021~\cite{cne_cp_1_2021}. Essa escolha garante plena aderência regulatória do curso perante o MEC/INEP, condição necessária para seu reconhecimento como Curso Superior de Tecnologia.

Dentro dessa denominação regulatória, o curso adota a ênfase curricular em \ppcenfasecurricular\ (\autoref{chp:identificacao}), que expressa a identidade pedagógica proposta neste PPC: além do núcleo tecnológico central de eletricidade, eletrônica, instrumentação, controle e automação — comum a qualquer curso do eixo Controle e Processos Industriais —, a matriz curricular deste curso incorpora um conjunto de conteúdos diferenciadores em robótica, sistemas inteligentes e inteligência artificial aplicada à automação (catalogados na aba \texttt{Conteudos} da matriz curricular sob o prefixo \texttt{AV\_}, "diferencial institucional aderente ao perfil de automação industrial"). Essa ênfase é comunicada institucionalmente pelo nome de divulgação \ppccursomercadologico, sem substituir a denominação regulatória em nenhum documento oficial (ver \autoref{chp:identificacao}).

\subsection{Justificativa da Organização em Área Básica de Ingresso (ABI)}

O \ppccurso é organizado na ABI juntamente com o curso de Engenharia de Automação, Robótica e Inteligência Artificial (\autoref{chp:abi}), compartilhando os cinco primeiros períodos da trajetória acadêmica. Essa articulação se justifica pela proximidade das áreas tecnológicas envolvidas — automação industrial, robótica e inteligência artificial aplicada — e pela oportunidade de oferecer ao estudante uma base formativa comum antes de exigir a opção entre uma formação tecnológica de ciclo curto e diretamente aplicada (o \ppccurso, \autoref{chp:justificativa}) e uma formação em Engenharia, de ciclo mais longo e com atribuições profissionais próprias.

A proposta de ABI busca ampliar a flexibilidade acadêmica sem descaracterizar as duas formações: o \ppccurso preserva identidade regulatória, matriz e perfil de egresso próprios (\autoref{chp:identificacao}; \autoref{chp:perfil_egresso}), assim como a Engenharia. A conclusão do \ppccurso resulta em formação tecnológica completa e terminal. Qualquer continuidade posterior na Engenharia sem novo processo seletivo é apenas uma possibilidade a ser instituída e regulamentada pela UFU, conforme o \autoref{sec:permanencia}.

\subsection{Vagas Ofertadas}

O \ppccurso oferta \textbf{\ppcvagasofertadas}, no turno \ppcturno, viabilizadas pelas novas vagas docentes e de técnicos de laboratório concedidas à FEELT no âmbito do Programa Universidades Inovadoras e Sustentáveis~\cite{programa_universidades_inovadoras_sustentaveis} (\autoref{sec:historico}). O dimensionamento dessa oferta inicial resulta da análise conjunta de quatro fatores: a capacidade acadêmica já instalada na FEELT; a densidade e a composição do parque industrial da região de influência de Uberlândia; a existência de uma lacuna efetiva de oferta pública, federal, presencial e noturna em Automação Industrial; e a compatibilidade entre a demanda potencial estimada e a capacidade de absorção profissional do mercado regional. Os elementos que fundamentam essa análise foram levantados no Estudo de Viabilidade para Criação de um Curso Superior de Tecnologia Noturno em Automação Industrial na UFU~\cite{ufu_estudo_viabilidade_cst_2026} e confrontados, nesta redação, com fontes institucionais primárias atualizadas.

\subsubsection{Contexto Regional e Ecossistema Produtivo}

A Regional Vale do Paranaíba da Federação das Indústrias do Estado de Minas Gerais (FIEMG)~\cite{fiemg_painel_vale_paranaiba_2026} -- que reúne 14 municípios sob a sede de Uberlândia, entre eles Araguari, Indianópolis, Monte Carmelo, Prata e Tupaciguara -- registrava, na edição de março de 2026 do Painel Regional da Indústria Mineira, 988.493 habitantes (Censo 2022), Produto Interno Bruto de R\$ 55,705 bilhões (2021, 6,5\% do PIB de Minas Gerais) e valor adicionado industrial de R\$ 13,416 bilhões (2021, 5,2\% do valor adicionado industrial mineiro), dos quais Uberlândia respondia por 78,8\%, seguida por Araguari (14\%) e Indianópolis (3,1\%). Em 2024, a indústria de transformação regional reunia 2.173 empresas e 43.714 empregos formais, e o setor de Serviços Industriais de Utilidade Pública (SIUP), 82 empresas e 2.366 empregos -- totalizando, com a indústria extrativa e a construção, 64.225 empregos industriais formais no recorte regional. O setor de alimentos concentrava isoladamente 31,3\% dos empregos industriais regionais, com apenas 11,0\% das empresas, o que evidencia forte presença de grandes plantas de processo contínuo. No acumulado de 12 meses até fevereiro de 2026, o saldo de emprego industrial da regional foi positivo, puxado por alimentos (+631 postos), bebidas (+186) e produtos de metal (+125), ao passo que manutenção e instalação de máquinas e equipamentos registrou saldo negativo (-248) -- sinal, segundo o próprio Estudo de Viabilidade, não de contração do setor, mas de reposição, rotatividade e necessidade de requalificação de mão de obra técnica~\cite{fiemg_painel_vale_paranaiba_2026,ufu_estudo_viabilidade_cst_2026}. Uberlândia é classificada pelo IBGE como Capital Regional B na Rede de Influência das Cidades (REGIC 2018)~\cite{ibge_regic_2018_uberlandia}, condição que reforça sua centralidade tanto urbana quanto produtiva no Triângulo Mineiro.

Esse parque industrial é concretamente aderente à automação, e não apenas em caráter genérico. Em alimentos, ingredientes e amidos, a Cargill mantém unidade em Uberlândia desde 1986 e investiu R\$ 150 milhões na ampliação de suas linhas de nutrição animal e amidos modificados, com geração de 600 empregos na obra e 20 vagas efetivas diretas~\cite{cargill_uberlandia_investimento}; em bebidas, a Ambev opera em Uberlândia a maior fábrica de bebidas do país, que recebeu novo aporte de R\$ 1,3 bilhão anunciado em dezembro de 2025 -- ampliando de três para cinco linhas de envase e dobrando a capacidade de marcas premium --, somando R\$ 2 bilhões investidos desde a inauguração em 2016 e mais de 66 mil empregos diretos, indiretos e induzidos gerados na cadeia produtiva~\cite{ambev_uberlandia_investimento_2025}; no complexo de proteínas e alimentos, a BRF mantém cinco unidades fabris em Uberlândia -- aves, óleo e gordura refinados, industrializados, suínos e margarinas --, com cerca de 290 produtores integrados, capacidade de 33 mil toneladas/mês e investimento de R\$ 319 milhões anunciado em 2021 para modernização e expansão~\cite{brf_uberlandia_unidade}; e, no corredor Indianópolis--Araguari, a LD Celulose -- \textit{joint venture} entre a austríaca Lenzing e a brasileira Dexco, com aporte de R\$ 6,3 bilhões e operação desde 2022 -- mantém planta de 500 mil toneladas anuais de celulose solúvel e 144 MW de energia limpa, tendo gerado mais de 10 mil empregos no Triângulo Mineiro~\cite{ld_celulose_planta}. Não por acaso, papel e celulose já respondem por 42,1\% das exportações industriais da regional, atrás apenas de alimentos (48,6\%)~\cite{fiemg_painel_vale_paranaiba_2026}. A esses somam-se as parcerias já mantidas por docentes da FEELT em pesquisa, desenvolvimento e inovação com ANEEL, CEMIG, Algar Telecom e ARCOM (\autoref{chp:justificativa}). Sem pretender esgotar o parque empresarial regional, esses exemplos demonstram a presença de indústrias de processo contínuo e de manufatura de alta instrumentação -- usuárias típicas de automação, supervisão, manutenção preditiva e integração industrial --, exatamente o campo de atuação profissional para o qual este curso forma.

A demanda por competências em automação também aparece no mercado de trabalho cotidiano: em abril de 2026, o Sistema Nacional de Emprego (SINE) de Uberlândia publicou vagas para técnico de manutenção industrial, eletricista de instalações industriais e operador de processo de produção, com remuneração entre aproximadamente R\$ 3,8 mil e R\$ 8 mil~\cite{sine_uberlandia_vagas_2026}.

\subsubsection{Lacuna de Oferta Formativa e Benchmark Nacional}

O levantamento da oferta concorrente mostra um mercado fragmentado, não saturado. A própria UFU oferta, na FEELT, o bacharelado em Engenharia de Controle e Automação -- integral, dez períodos, 30 vagas autorizadas, das quais 15 pelo SISU, até a presente reformulação (\autoref{chp:justificativa} do PPC de Engenharia) --; o Instituto Federal do Triângulo Mineiro (IFTM) mantém Tecnologia em Automação Industrial, porém no \textit{campus} de Ituiutaba, fora de Uberlândia~\cite{iftm_automacao_industrial}; e o SENAI de Uberlândia concentra oferta de cursos técnicos e de aperfeiçoamento em CLP, instrumentação, redes industriais, eletricista industrial e inversores, sem ofertar graduação~\cite{ufu_estudo_viabilidade_cst_2026}. A oferta noturna em engenharias correlatas hoje existente em Uberlândia (Engenharia de Controle e Automação e Engenharia Elétrica) é de natureza privada, mantida pela Uniube, com dez semestres de duração~\cite{ufu_estudo_viabilidade_cst_2026}. Não foi identificada, portanto, oferta pública, federal, presencial e noturna de graduação tecnológica em Automação Industrial em Uberlândia -- lacuna que este curso preenche, situando-se no espaço formativo intermediário entre a formação técnica de nível médio e a Engenharia, e não em competição direta com nenhuma delas.

Esse desenho -- tecnólogo noturno, aplicado, industrial, de três anos -- é praticado em instituições públicas federais de referência. A Universidade Tecnológica Federal do Paraná (UTFPR) mantém, em Ponta Grossa, Tecnologia em Automação Industrial presencial, noturna, com nota 5 na avaliação do MEC~\cite{utfpr_automacao_industrial}; o Instituto Federal de São Paulo (IFSP) oferece o mesmo eixo tecnológico em múltiplos \textit{campi}; e a rede de Faculdades de Tecnologia (Fatec) do Centro Paula Souza oferta Automação Industrial e Mecatrônica Industrial, inclusive no turno noturno, com currículo aplicado a PLC, robótica e manufatura~\cite{ufu_estudo_viabilidade_cst_2026}. A UTFPR mantém, ainda, tabela formal de equivalência entre disciplinas do tecnólogo em Automação Industrial e da engenharia correspondente~\cite{utfpr_equivalencia_automacao_eng}, evidenciando que os dois percursos podem ser tratados institucionalmente como trajetórias complementares -- não concorrentes --, o que reforça a coerência da articulação entre este curso e a Engenharia de Automação, Robótica e Inteligência Artificial na Área Básica de Ingresso (\autoref{chp:abi}).

\subsubsection{Dimensionamento das 30 Vagas: Demanda e Capacidade Institucional}

Do lado da demanda, o Estudo de Viabilidade estimou o público potencial a partir do estoque de trabalho industrial formal da região -- trabalhadores CLT da indústria, técnicos que buscam verticalização, egressos do SENAI e do IFTM, e estudantes-trabalhadores hoje sem acesso a uma graduação noturna equivalente. Considerando o estoque de 64.225 empregos industriais formais do recorte regional, taxas de conversão conservadoras de 0,25\% a 0,60\% desse estoque em candidatos anuais projetam entre aproximadamente 161 e 385 candidatos potenciais por ano~\cite{ufu_estudo_viabilidade_cst_2026} -- valor amplamente superior às 30 vagas ofertadas, indicando escala inicial prudente, e não expansionista.

Do lado da absorção profissional, o mesmo estudo aplicou critério deliberadamente conservador, restrito às 2.173 empresas de transformação e às 82 de SIUP (2.255 estabelecimentos), sem incluir construção, logística fora da CNAE industrial ou serviços técnicos externos. Nesse recorte, cenários de 2,0\%, 3,5\% e 5,0\% de estabelecimentos absorvendo ou requalificando ao menos um profissional de perfil tecnólogo por ano projetam, respectivamente, 45, 79 e 113 profissionais/ano -- e 225, 395 e 565 em cinco anos~\cite{ufu_estudo_viabilidade_cst_2026}. Combinando as duas pontas, o estudo concluiu que uma coorte anual entre 30 e 40 vagas está em escala compatível tanto com a demanda potencial quanto com a capacidade de absorção regional. A oferta de 30 vagas adotada neste PPC situa-se no extremo prudencial dessa faixa recomendada, e não em seu limite superior -- opção institucionalmente conservadora, que privilegia a consolidação da oferta antes de eventual ampliação futura.

Do lado da capacidade interna, a oferta apoia-se nas novas vagas docentes e de técnicos de laboratório concedidas à FEELT pelo Programa Universidades Inovadoras e Sustentáveis~\cite{programa_universidades_inovadoras_sustentaveis}, somadas à estrutura já consolidada da unidade -- 1.373 estudantes, 72 docentes e 16 laboratórios em funcionamento, dos quais 10 docentes vinculados ao Departamento de Sistemas de Controle e Automação e o Laboratório de Automação e Sistemas Elétricos de Controle (LASEC), com atuação já estabelecida em automação industrial, redes industriais, sistemas embarcados, robótica, máquinas elétricas e acionamentos, e interação com empresas da região~\cite{feelt_numeros}. \textcolor{red}{[INSERIR NÚMERO DE VAGAS DOCENTES EFETIVAMENTE CONCEDIDAS PELO PROGRAMA UNIVERSIDADES INOVADORAS E SUSTENTÁVEIS, POR ÁREA/DEPARTAMENTO; INSERIR NÚMERO DE VAGAS DE TÉCNICOS DE LABORATÓRIO EFETIVAMENTE CONCEDIDAS; CONFIRMAR JUNTO AO NDE/COLEGIADO A CAPACIDADE NOMINAL DAS BANCADAS E LABORATÓRIOS DESTINADOS ÀS ATIVIDADES PRÁTICAS DESTE CURSO E A EVENTUAL NECESSIDADE DE NOVAS AQUISIÇÕES DE EQUIPAMENTOS -- DADOS NÃO LOCALIZADOS NOS DOCUMENTOS DISPONÍVEIS PARA ESTE PPC]} A alocação de 30 vagas -- e não das até 40 vagas admitidas pelo próprio Estudo de Viabilidade -- reflete, portanto, o equilíbrio entre uma demanda regional que comportaria turmas maiores e a necessidade de preservar relação adequada entre docentes, técnicos e estudantes nas atividades práticas de laboratório, condição essencial de segurança e de qualidade pedagógica em um curso de natureza eminentemente aplicada.

\subsubsection{O Turno Noturno como Decisão Estratégica}

A oferta no turno noturno não é uma característica meramente operacional deste curso: é uma decisão acadêmica, social e regional deliberada. Ela viabiliza o ingresso e a permanência de um público que, de outro modo, permaneceria estruturalmente afastado do ensino superior público federal presencial -- trabalhadores já empregados na indústria regional, técnicos em automação, eletrotécnica, eletrônica e mecatrônica, eletricistas industriais, instrumentistas, operadores de processo e egressos de cursos técnicos do SENAI, do IFTM e de instituições correlatas que hoje buscam verticalização acadêmica sem interromper seu vínculo profissional. O bacharelado em Engenharia de Controle e Automação da própria UFU sempre operou em regime integral, com barreira de horário estrutural para o público trabalhador~\cite{ufu_estudo_viabilidade_cst_2026}; e a oferta noturna hoje existente em áreas correlatas em Uberlândia -- Engenharia de Controle e Automação e Engenharia Elétrica, ambas na Uniube, com dez semestres de duração -- é de natureza privada e mais onerosa para o estudante~\cite{ufu_estudo_viabilidade_cst_2026}. O turno noturno deste curso preenche, portanto, uma lacuna concreta -- e não meramente hipotética -- de acesso público federal gratuito no período noturno em Uberlândia.

\subsubsection{Integração com Estágio, Indústria e Formação Aplicada}

A permanência do(a) estudante no ambiente industrial durante o dia constitui, para este curso, uma vantagem pedagógica, e não apenas uma condição a ser acomodada. Ela favorece a identificação de problemas reais da indústria, o estágio e a atividade profissional correlata à formação, o desenvolvimento de projetos aplicados, de trabalhos de conclusão de curso e de projetos integradores, e a aproximação contínua entre a universidade e o setor produtivo -- inclusive por meio das Atividades Curriculares de Extensão previstas pela Resolução CONGRAD nº 177/2026~\cite{ufu_congrad_177_2026} (\autoref{sec:ace}). Configura-se, assim, um potencial ciclo virtuoso entre indústria, universidade e estudante: um problema tecnológico real identificado no ambiente de trabalho pode tornar-se objeto de projeto acadêmico, cuja solução, uma vez desenvolvida, retorna como competência aplicada ao próprio ambiente produtivo. Estudantes que já atuam profissionalmente em ambientes industriais podem, nesse sentido, trazer para as disciplinas experiências, problemas e demandas tecnológicas concretas que enriquecem discussões, projetos e atividades práticas -- sempre respeitados a confidencialidade, a propriedade intelectual das empresas envolvidas e os regulamentos acadêmicos da UFU.

A composição do corpo docente do curso pode se beneficiar dessa proximidade com o setor produtivo de forma direta e institucionalmente sólida. A Lei nº 12.772/2012, no art. 20, estabelece, ao lado do regime de 40 horas semanais com Dedicação Exclusiva (DE) -- que veda, salvo exceções pontuais e limitadas, o exercício de outra atividade remunerada --, o regime de tempo parcial de 20 horas semanais, sem essa vedação~\cite{lei_12772_2012}; a UFU regulamenta o ingresso e a alteração de regime de trabalho docente por meio da PROGEP~\cite{ufu_progep_regime_trabalho}. À luz desse marco legal, é interesse do curso compor parte de seu quadro docente efetivo -- via concurso público regular, observados os limites orçamentários e de Banco de Professor-Equivalente aplicáveis -- também com docentes em regime de 20 horas, e não exclusivamente em Dedicação Exclusiva. O objetivo é criar, dentro do próprio quadro docente, um vínculo estável com profissionais da região que permanecem em atividade no mercado de trabalho durante o dia e lecionam à noite -- trazendo para a sala de aula experiências, processos e problemas tecnológicos vividos concretamente na indústria, contribuindo com as Atividades Curriculares de Extensão do curso~\cite{ufu_congrad_177_2026} (\autoref{sec:ace}) e aproximando os(as) estudantes das empresas da região. Esse mecanismo é distinto -- e não deve ser confundido com -- a contratação de professor substituto ou visitante, de natureza temporária e finalidade diversa (Resolução CONSUN nº 24/1994)~\cite{ufu_consun_24_1994}. \textcolor{red}{[DEFINIR, JUNTO AO NDE E À DIREÇÃO DA FEELT, A PROPORÇÃO OU O NÚMERO DE VAGAS DOCENTES DO CURSO A SEREM DESTINADAS A REGIME DE 20 HORAS DENTRO DO QUANTITATIVO TOTAL DE VAGAS CONCEDIDAS PELO PROGRAMA UNIVERSIDADES INOVADORAS E SUSTENTÁVEIS]}

\subsubsection{Diferenciação em Relação à Engenharia e Articulação entre os Dois Cursos}

Este curso não deve ser compreendido como uma "engenharia reduzida", mas como uma graduação tecnológica de identidade própria, com finalidade profissional específica: formação aplicada, voltada à integração, implantação, operação e manutenção de sistemas automatizados, ao \textit{troubleshooting}, e ao domínio operacional de CLP, SCADA, redes industriais, instrumentação e acionamentos elétricos, em conexão direta com as demandas operacionais e tecnológicas da indústria regional. A Engenharia de Automação, Robótica e Inteligência Artificial, por sua vez, oferece formação científica e tecnológica mais ampla, voltada à modelagem, análise, projeto e desenvolvimento de sistemas, ao controle avançado, à robótica, aos sistemas embarcados e à inteligência artificial aplicada, à pesquisa e à inovação, e à atuação em engenharia e liderança técnica (\autoref{chp:justificativa} do respectivo PPC). Os dois perfis são complementares, não hierárquicos -- distinção que o próprio benchmark nacional confirma: o bacharelado da UFU antecedente a esta reformulação já operava com 30 vagas autorizadas, das quais 15 pelo SISU, em regime integral, o que demonstra que os dois formatos sempre coexistiram como trajetórias paralelas, não como versões concorrentes de uma mesma formação.

Essa complementaridade é reforçada pela organização conjunta dos dois cursos em Área Básica de Ingresso (\autoref{chp:abi}), que compartilham os cinco primeiros períodos da trajetória acadêmica. A implantação articulada permite racionalizar a infraestrutura laboratorial -- em particular sua utilização no período noturno --, compartilhar de forma planejada componentes curriculares básicos e tecnológicos comuns, fortalecer os grupos de pesquisa e extensão da área de Controle e Automação da FEELT, e ampliar a interação conjunta com as empresas da região. O objetivo central dessa articulação, contudo, não é a economia de recursos, mas a sinergia acadêmica e a utilização responsável da capacidade institucional já existente na FEELT.

\subsubsection{Transformação Digital da Indústria}

A relevância desta oferta não decorre apenas do parque industrial hoje instalado, mas também de sua transformação tecnológica em curso. No âmbito da Nova Indústria Brasil, o Ministério do Desenvolvimento, Indústria, Comércio e Serviços (MDIC) estabeleceu a meta de digitalizar 50\% das empresas industriais brasileiras até 2033, com meta intermediária de 25\% até 2026 -- partindo de 18,9\% em 2023~\cite{mdic_nova_industria_brasil_2024} --, e relançou o programa Brasil Mais Produtivo, com R\$ 2 bilhões destinados a atender 200 mil empresas até 2027, em parceria com BNDES, Finep, Embrapii, ABDI, Sebrae e SENAI~\cite{mdic_nova_industria_brasil_2024}. A Federação Internacional de Robótica registrou 542 mil robôs industriais instalados globalmente em 2024, mais do que o dobro de dez anos antes, e a ABDI associa a evolução do controle e automação industrial à Internet Industrial das Coisas (IIoT) e à manufatura avançada~\cite{abdi_boas_praticas_produtivas_2024}. Essas tendências reforçam, e não reduzem, a relevância de uma formação tecnológica aplicada em automação industrial. A matriz curricular deste curso incorpora esses eixos por meio de conteúdos em PLC/SCADA, redes industriais, sensores e dados industriais, manutenção preditiva, integração OT/IT, visão computacional aplicada e fundamentos de inteligência artificial industrial (\autoref{chp:estrutura_curricular}).

\subsubsection{Acompanhamento e Revisão Futura}

A dimensão inicial de 30 vagas representa uma coorte tecnicamente compatível com a capacidade institucional hoje instalada na FEELT e com a demanda regional estimada neste PPC -- e não um limite definitivo. O Núcleo Docente Estruturante e o Colegiado do Curso deverão acompanhar periodicamente indicadores de procura por vaga, ocupação, retenção, evasão, progressão acadêmica, inserção em estágio, empregabilidade e atuação profissional dos egressos, e utilização da infraestrutura laboratorial, de modo a subsidiar eventual revisão do quantitativo de vagas com base em evidências, e não em projeções não fundamentadas.

\section{Histórico do Curso}
\label{sec:historico}

O \ppccurso encontra-se, no momento da elaboração deste PPC, em fase de proposta (\texttt{perfil.status: proposta}, currículo \ppcversao), com Portaria de Autorização em tramitação no MEC (\autoref{chp:identificacao}). Trata-se de uma proposta de curso inteiramente novo — e não de uma reformulação ou renomeação de curso de graduação já existente na FEELT —, distinta nesse sentido da Engenharia de Automação, Robótica e Inteligência Artificial, que é reformulação e renomeação do curso já existente de Engenharia de Controle e Automação (\autoref{chp:abi}). Os dois cursos, ambos contemplados pelo Programa Universidades Inovadoras e Sustentáveis~\cite{programa_universidades_inovadoras_sustentaveis} — este pela criação, aquele pela reformulação —, foram organizados conjuntamente em Área Básica de Ingresso (ABI) em razão dessa origem institucional comum.

% REVISAR: histórico formal de criação do curso (ato/resolução de criação, data de aprovação pelos conselhos superiores da UFU, ano previsto de ingresso da primeira turma) não localizado nos documentos ou na matriz curricular deste perfil — a aba \texttt{Legislacao} registra a base normativa do CNCST, das DCN EPT e do Programa Universidades Inovadoras e Sustentáveis, mas não um ato de criação institucional específico deste curso, nem o edital/portaria exato de contemplação pelo programa.

\section{Unidade Acadêmica}

O \ppccurso está vinculado à Faculdade de Engenharia Elétrica (FEELT) da Universidade Federal de Uberlândia (UFU), unidade acadêmica responsável por sua gestão pedagógica, administrativa e acadêmica. A FEELT atualmente abriga seis cursos de graduação presenciais — cinco localizados no campus Santa Mônica, em Uberlândia/MG (Engenharia Elétrica, Engenharia Biomédica, Engenharia de Computação, Engenharia Eletrônica e de Telecomunicações, Engenharia de Automação, Robótica e Inteligência Artificial) e um no campus Patos de Minas/MG (Engenharia Eletrônica e de Telecomunicações) — além de dois programas de pós-graduação \textit{stricto sensu}: o Programa de Pós-Graduação em Engenharia Elétrica e o Programa de Pós-Graduação em Engenharia Biomédica.

Com a aprovação do Regimento Interno da FEELT~\cite{ufu_feelt_regimento_2023}, a estrutura organizacional da unidade está reorganizada em Departamentos Acadêmicos, cada um responsável por atividades de ensino, pesquisa e extensão em áreas específicas do conhecimento. Os departamentos atualmente constituídos na FEELT são:

\begin{itemize}
    \item Departamento de Sistemas de Potência e Energia (DSPE)
    \item Departamento de Engenharia Biomédica (DEB)
    \item Departamento de Computação e Inovação em Engenharia (DCIE)
    \item Departamento de Sistemas de Controle e Automação (DSCA)
    \item Departamento de Eletrônica e Telecomunicações (DETEL)
    \item Departamento de Máquinas Elétricas e Materiais (DMEM)
\end{itemize}

Pela natureza de sua matriz curricular — concentrada em eletricidade, eletrônica, instrumentação, controle, automação e redes industriais (\autoref{chp:identificacao}; aba \texttt{Nucleos} da matriz curricular) —, o \ppccurso tem forte interface com o Departamento de Sistemas de Controle e Automação (DSCA), com contribuição de docentes de outros departamentos da FEELT (DETEL, DMEM, DCIE) e de unidades acadêmicas externas, como a Faculdade de Matemática (FAMAT), o Instituto de Física (INFIS) e a Faculdade de Engenharia Química (FEQUI), na oferta de disciplinas obrigatórias (aba \texttt{Componentes} da matriz curricular). % REVISAR: unidade/departamento formalmente designado como responsável administrativo pelo curso ainda não confirmado nos documentos deste perfil.

Essa estrutura acadêmica consolidada garante ao curso uma base sólida de infraestrutura laboratorial e corpo docente qualificado para a formação de tecnólogos aptos a atuar nos mais diversos contextos da automação industrial.

\section{Relação entre Sociedade e o Curso}

O \ppccurso orienta-se pela missão de formar profissionais capazes de responder às necessidades do setor produtivo, contribuindo para o avanço tecnológico, econômico e social da região e do país. Sua proposta articula-se com o contexto local e regional, buscando integrar-se ao ecossistema de inovação do Triângulo Mineiro.

A cidade de Uberlândia/MG, onde o curso está localizado, destaca-se como um dos principais polos tecnológicos e industriais do interior do Brasil, combinando infraestrutura logística estratégica, mão de obra qualificada e ambiente favorável ao desenvolvimento de negócios de base tecnológica. Diversas instituições e iniciativas locais vêm consolidando esse ecossistema, como a Secretaria Municipal de Desenvolvimento Econômico, Inovação e Turismo, a i9 Uberlândia, a Comunidade Colmeia, a Minas Startup, o Parque do Sabiá Tech e o SEBRAE/MG, fomentando a articulação entre universidades, indústrias, startups e o poder público.

A atuação de docentes da FEELT em projetos de pesquisa, desenvolvimento e inovação (PD\&I) tem resultado em parcerias com empresas e instituições dos setores elétrico e industrial — como ANEEL, CEMIG, Algar Telecom e ARCOM — e em projetos financiados por agências de fomento nacionais (CAPES, CNPq, FINEP, FAPEMIG), diretamente relacionadas ao campo de atuação de um(a) Tecnólogo(a) em Automação Industrial. % REVISAR: parcerias, convênios e indicadores de inserção profissional específicos deste curso (ainda em fase de proposta, sem egressos) não estão disponíveis; os dados acima refletem a atuação institucional da FEELT como um todo, não deste curso especificamente.

A proposta deste PPC reafirma o compromisso da UFU com a sociedade, respondendo à demanda por educação pública de qualidade voltada para uma área de alta relevância tecnológica e industrial, e promovendo uma formação alinhada às exigências contemporâneas da Indústria 4.0, da sustentabilidade e da inovação.
